% Part: proof-theory
% Chapter: natural-deduction
% Section: sequents

\documentclass[../../../include/open-logic-section]{subfiles}

\begin{document}

\olfileid{pt}{nat}{seq}

\olsection{له سېکوېنټونو سره فطري استنتاج: \Log{N2c} او~\Log{N2i}}

په \Log{N1c} يا \Log{N1i} کښې !!^a{proof}~$\delta$
ښيي چې پاييز !!{formula} ئې~$!A$ له پرانيستو فرضونو څخه
نتيجه کېږي۔ که $\Gamma$ د هغو !!{formula}s~$!B$ سټ وي چې
$!B^x$ ئې په~$\delta$ کښې د پرانيستي فرض په توګه راځي، نو
دا په $\delta \Proves \Gamma
\Sequent !A$ ليکلے شو۔ دا ښيي چې د سېکوېنټونو دپاره هم يو
فطري استنتاجي نظام جوړولے شو۔ په داسې نظام کښې هر سېکوېنټ دا
معلومات لري چې ښي اړخ ته !!{formula} په کومو پرانيستو فرضونو
پورې تړلے دے، يعنې کوم فرضونه په هغه فرعي !!{proof} کښې
پرانيستي دي چې په دغه سېکوېنټ ختمېږي۔ دا له سېکوېنټونو سره
فطري استنتاج، يا په ``سېکوېنټي طرز'' فطري استنتاج دے؛ پايله
شوي نظامونه \Log{N2c} او~\Log{N2i} بولو۔
\textbf{سرچينه‌يي سمون OLCMP-145:} د دې لومړني ثبوت ونې
د N1 دي؛ د سرچينې نقل شوي G1 نومونه د همدې برخې له پرتله او
دواړو ورپسې اصلي قضيو سره سم شول۔

د \Log{N1} او \Log{N2} د لا نېغ سمون دپاره، کيڼ اړخ ته
سټونه~$\Gamma$ د نښه لرونکو !!{formula}s~$x:!B$ څخه جوړ
اخلو۔ په \Log{N1} کښې قاعدې يوازې د مقدمو پر !!{formula}s
کار کوي، يعنې د سمدستي فرعي !!{proof}s پر پاييزو
!!{formula}s؛ نو د \Log{N2} قاعدې يوازې د سېکوېنټ پر ښي
فارمول کار کوي۔ په دې توګه د سېکوېنټونو کيڼ اړخ ته نښه لرونکي
فارمولونه يوازې سياق بللے شو۔

د فرضونو پر ځاے، د \Log{N2} ثبوتونو تر ټولو بره سېکوېنټونه
د دې بڼې بديهيات دي:
\[x:!A \Sequent !A.\]
قاعدې د \Log{N1} له قواعدو سره کره سمون لري او په
\olref{tab:N2} کښې ورکړل شوې دي۔

\subfile{rules-N2}

لکه په \Log{N1c} او \Log{N1i} کښې، بايد دا ومنو چې
\Intro{\lif}، \Intro{\lnot}، \Elim{\lor} او \Elim{\lexists}
قاعدې فرضونه ايستلے شي، خو اړې نۀ دي چې فرضونه وباسي۔ نو په
\Log{N2c} او~\Log{N2i} کښې د اړوندو قواعدو سمې کارونې هغه
وخت هم منو چې اړوند سياقي !!{formula}s $!A$، $!B$ او
$!A(a)$ د اړوندې مقدمې په سياق کښې نۀ وي۔

\begin{prop}
که $\delta$ په~\Log{N1c} يا \Log{N1i} کښې د~$!A$
!!a{proof} وي، نو په~\Log{N2c} (\Log{N2i}) کښې د
$\Gamma \Sequent !A$ !!a{proof}~$\delta'$ شته، چې
$\Gamma$ د ټولو هغو~$x:!B$ سټ دے چې $!B^x$ ئې د~$\delta$
پرانيستے فرض وي۔
\end{prop}

\begin{proof}
  په $n=\pheight{\delta}$ استقرا کوو۔ که $n=0$ وي، $\delta$
  يوازې يو پرانيستے فرض~$!A^x$ دے۔ نو $\Gamma = \{x:!A\}$
  او $\Gamma
  \Sequent !A$ د \Log{N2c} يا~\Log{N2i} يو بديهيت دے۔

  که $n>0$ وي، د~$\delta$ د وروستي استنتاج له مخې حالتونه
  بېلوو۔ يوازې يو حالت بيان کوو؛ پاتې حالتونه تمرينونه دي۔

  فرض کړئ وروستے استنتاج $\Intro{\lif}$ وي۔ نو $\delta$
  په دې ختمېږي:
  \[
  \AxiomC{$\Discharge{!B}{x}$}
  \RightLabel{$\delta_1$}
  \DeduceC{$!C$}
  \DischargeRule{\Intro{\lif}}{x}
  \UnaryInfC{$!B \lif !C$}
  \DisplayProof
  \]
  د استقرا له فرضيې، په~\Log{N2c} (\Log{N2i}) کښې د
  $\Gamma' \Sequent !C$ !!a{proof}~$\delta_1'$ شته، چې
  $\Gamma'$ هغه نښه لرونکي فرضونه دي چې په~$\delta_1$
  کښې پرانيستي وي۔ که $!B^x$ د $\delta_1$ پرانيستے فرض وي،
  نو د \Intro{\lif} په استنتاج ايستل کېږي او د~$\delta$
  پرانيستي فرضونه $\Gamma =
  \Gamma' \setminus \{x:!B\}$ دي۔ يعنې $\delta_1'$ د
  $x:!B, \Gamma \Sequent !C$ !!a{proof} دے؛ نو $\delta'$
  داسې اخستلے شو:
  \[
  \AxiomC{}
  \RightLabel{$\delta_1'$}
  \Deduce$x:!B, \Gamma \fCenter !C$
  \RightLabel{\Intro{\lif}}
  \UnaryInf$\Gamma \fCenter !B \lif !C$
  \DisplayProof
  \]
  \textbf{سرچينه‌يي سمون OLCMP-146:} د رسم په سياق کښې
  هماغه شېمايي فارمول پکار دے؛ د سرچينې پاتې شوې نښه ورزياته شوه۔
  که $!B^x$ د~$\delta_1$ پرانيستے فرض نۀ وي، نو
  \Intro{\lif} هېڅ فرض نۀ باسي او د $\delta$ او $\delta_1$
  پرانيستي فرضونه يو شان دي۔ په~\Log{N2c} او~\Log{N2i}
  کښې د \Intro{\lif} د مقدمې په سياق کښې د $x:!B$ شته‌والے
  اړين نۀ دے؛ نو $\delta'$ داسې اخستلے شو:
  \[
  \AxiomC{}
  \RightLabel{$\delta_1'$}
  \Deduce$\Gamma \fCenter !C$
  \RightLabel{\Intro{\lif}}
  \UnaryInf$\Gamma \fCenter !B \lif !C$
  \DisplayProof
  \]

  هغه حالتونه چې $\delta$ په بله قاعده ختمېږي، همداسې دي۔
\end{proof}

\begin{prop}
که $\delta$ په \Log{N2c} يا \Log{N2i} کښې د $\Gamma
\Sequent !A$ !!a{proof} وي، نو په \Log{N1c} (\Log{N1i})
کښې !!a{proof}~$\delta'$ شته چې پاييز !!{formula} ئې~$!A$
وي او د هر $x:!B \in \Gamma$ دپاره پرانيستي فرضونه ئې
$!B^x$ وي۔
\end{prop}

\begin{proof}
  بيا د~$\delta$ پر !!{height}~$n$ استقرا کوو۔ که $n=0$
  وي، $\delta$ يو بديهيت دے، $x:!A \Sequent !A$۔
  !!{proof}~$\delta'$ هماغه فرض~$!A^x$ دے۔

  که $n>0$ وي، د~$\delta$ د وروستي استنتاج له مخې حالتونه
  بېلوو۔ مثلاً، که دغه استنتاج~$\Intro{\lif}$ وي، نو
  $\delta$~په دې ختمېږي:
  \[
  \bottomAlignProof
  \AxiomC{}
  \RightLabel{$\delta_1$}
  \Deduce$x:!B, \Gamma \fCenter !C$
  \RightLabel{\Intro{\lif}}
  \UnaryInf$\Gamma \fCenter !B \lif !C$
  \DisplayProof
  \quad\text{ يا }\quad
  \bottomAlignProof
  \AxiomC{}
  \RightLabel{$\delta_1$}
  \Deduce$\Gamma \fCenter !C$
  \RightLabel{\Intro{\lif}}
  \UnaryInf$\Gamma \fCenter !B \lif !C$
  \DisplayProof
  \]
  د استقرا له فرضيې، په~\Log{N1c} (\Log{N1i}) کښې د~$!C$
  !!a{proof}~$\delta_1'$ شته۔ په لومړي حالت کښې $!B^x$
  د~$\delta_1'$ پرانيستے فرض دے؛ په دويم کښې نۀ دے۔
  !!{proof}~$\delta'$ په ترتيب داسې اخستلے شو:
  \[
  \bottomAlignProof
  \AxiomC{$\Discharge{!B}{x}$}
  \RightLabel{$\delta_1'$}
  \DeduceC{$!C$}
  \DischargeRule{\Intro{\lif}}{x}
  \UnaryInfC{$!B \lif !C$}
  \DisplayProof
  \quad\text{ يا }\quad
  \bottomAlignProof
  \AxiomC{}
  \RightLabel{$\delta_1'$}
  \DeduceC{$!C$}
  \RightLabel{\Intro{\lif}}
  \UnaryInfC{$!B \lif !C$}
  \DisplayProof
  \]
  \textbf{سرچينه‌يي سمون OLCMP-147:} د پرانيستي فرض خبره
  د استقرا له مخې جوړ شوي سمدستي فرعي ثبوت ده، نۀ د پاييز
  ثبوت چې فرض ترې ايستل شوے دے۔ په وروستي رسم کښې هم د همدې
  نوي فرعي ثبوت نښه سمه شوه۔
\end{proof}

\end{document}
